\documentclass[../../../include/open-logic-section]{subfiles}

\begin{document}

\olfileid{sth}{story}{urelements}
\olsection{اورعناصر یا نه؟}

در چند فصل آینده خواهیم کوشید اصول موضوع را از تصور انباشتی-تکراریِ
مجموعه استخراج کنیم. اما پیش از آنکه جلوتر برویم، باید دربارهٔ
\emph{اورعناصر} اندکی بیشتر سخن بگوییم.

تصویرِ~\olref[approach]{sec} تنها به ما اجازه می‌داد مجموعه‌های تازه را
از \emph{مجموعه‌های} پیشین بسازیم. بااین‌حال، شاید بخواهیم اجازه دهیم
برخی \emph{نامجموعه‌ها}---گاوها، خوک‌ها، دانه‌های شن یا هر چیز
دیگر---!!{element} مجموعه‌ها باشند. در این صورت، با برخی عناصر بنیادین،
یعنی \emph{اورعناصر}، آغاز می‌کنیم و سپس می‌گوییم در هر مرحلهٔ~$S$،
«همهٔ مجموعه‌های ممکن» متشکل از اورعناصر یا مجموعه‌های ساخته‌شده در مراحل
پیش از~$S$ (با هر ترکیبی از آن‌ها) را می‌سازیم. تصویر حاصل بیشتر به شکل
زیر خواهد بود:
\begin{center}
	\begin{tikzpicture}[scale=0.6]
	\tikzset{cut_here/.style={densely dotted}}
	\draw (-4,6) -- (-1, 0) -- (1,0) -- (4,6);
	\draw[cut_here] (-5,8)--(-4,6);
	\draw[cut_here] (5,8)--(4,6);
	\draw(-1.5, 1)--(1.5, 1);
	\draw(-2, 2)--(2, 2);
	\draw(-2.5, 3)--(2.5,3);
	\draw(-3, 4)--(3, 4);
	\draw(-3.5, 5)--(3.5, 5);
	\draw(-4, 6)--(4, 6);
	\node[label] at (2, 0) {\small 0};
	\node[label] at (2.5, 1) {\small 1};
	\node[label] at (3, 2) {\small 2};
	\node[label] at (3.5, 3) {\small 3};
	\node[label] at (4, 4) {\small 4};
	\node[label] at (4.5, 5) {\small 5};
	\node[label] at (5, 6) {\small 6};
	\end{tikzpicture}
\end{center}
اکنون باید تصمیم بگیریم: \emph{آیا باید اورعناصر را مجاز بدانیم؟}

از دیدگاه فلسفی، گنجاندن اورعناصر در نظریه‌پردازی ما معقول است. دلیل اصلی
آن است که نظریهٔ مجموعه‌هایمان را \emph{کاربردپذیر} سازیم. برای توضیح این
نکته، از~\olref[sfr][siz][]{chap} به یاد آورید که می‌گوییم دو مجموعهٔ
$A$ و~$B$ هم‌اندازه‌اند، یعنی~$\cardeq{A}{B}$، اگر و تنها اگر تناظری
یک‌به‌یک میان آن‌ها وجود داشته باشد. حال اگر گاوهای چراگاه و خوک‌های آغل
هر دو مجموعه‌ای تشکیل دهند، می‌توانیم گزارهٔ «تعداد گاوها با تعداد خوک‌ها
برابر است» را به زبان نظریهٔ مجموعه‌ها تحلیل کنیم. اما اگر اورعناصر را
ممنوع کنیم، به‌طوری‌که گاوها و خوک‌ها مجموعه تشکیل \emph{ندهند}، این
تحلیل نظریه‌مجموعه‌ای در دسترس نخواهد بود. در واقع، دیگر راه سرراستی برای
کاربرد نظریهٔ مجموعه‌ها بر چیزی جز خود مجموعه‌ها نخواهیم داشت. (برای
دلایل بیشتر در تأیید گنجاندن اورعناصر، بنگرید به
\citealt[pp.~vi, 24, 50--1]{Potter2004}.)

بااین‌همه، در ریاضیات مجازدانستن اورعناصر بسیار نادر است. یک دلیل آن است
که صورت‌بندی نظریهٔ مجموعه‌ها بدون اورعناصر \emph{اندکی} آسان‌تر است. اما
گاه توجیه‌های جالب‌تری برای کنارگذاشتن اورعناصر از نظریهٔ مجموعه‌ها می‌یابیم:
\begin{quote}
	مطابق با این باور که نظریهٔ مجموعه‌ها بنیاد ریاضیات است، باید بتوانیم
	تمام ریاضیات را تنها با سخن‌گفتن از مجموعه‌ها دربر بگیریم؛ پس متغیرهای
	ما نباید بر اشیایی مانند گاوها و خوک‌ها دامنه داشته باشند.
	%اما اگر $C$ یک گاو باشد، $\{C\}$ یک مجموعه است، ولی شیئی ریاضیِ مجاز نیست.
	\citep[p.~8]{Kunen1980}
\end{quote}
پس تمرکز بر کاربردپذیری، \emph{گنجاندن} اورعناصر را پیشنهاد می‌کند؛ تمرکز
بر هدفی بنیادگرایانه و تحویلی (فروکاستن ریاضیات به نظریهٔ مجموعه‌های محض)
ممکن است \emph{کنارگذاشتن} آن‌ها را پیشنهاد کند. اندکی تن‌آسایی نیز ما را
به سوی کنارگذاشتن اورعناصر سوق می‌دهد.

ما تن‌آسانه‌ترین راه را در پیش خواهیم گرفت. البته این انتخاب تا اندازه‌ای
توجیهی آموزشی نیز دارد. هدف ما آشناکردن شما با آن بخش‌هایی از نظریهٔ
مجموعه‌هاست که برای آغاز مطالعهٔ فلسفهٔ نظریهٔ مجموعه‌ها لازم‌اند. و بیشتر
آثار فلسفی این حوزه نظریه‌های مجموعه‌ای را بررسی می‌کنند که \emph{بدون}
اورعناصر صورت‌بندی شده‌اند. بنابراین، اگر این کتاب نسبتاً نزدیک به آن آثار
حرکت کند، شاید سودمندتر باشد.

\end{document}
